\documentclass[11pt,a4paper]{article}
\usepackage[utf8]{inputenc}
\usepackage{amsmath,amssymb,amsthm}
\usepackage{geometry}

\geometry{margin=1in}

\title{\textbf{Fuss--Catalan Series for the Algebraic Equation
        $x-bx^n=a$}}
\author{\textbf{Dawood Ali Hassoun}\\
\textit{Independent Mathematical Researcher, Latakia, Syria}}
\date{November 2026}

\newtheorem{theorem}{Theorem}
\newtheorem{lemma}{Lemma}
\newtheorem{remark}{Remark}

\begin{document}

\maketitle

\begin{abstract}
For a fixed integer $n\geq 2$, we derive the power-series branch through
$x=a$ of the equation $x-bx^n=a$. Lagrange--B\"urmann inversion gives the
Fuss--Catalan coefficients
\[
C_k^{(n)}=\frac{1}{nk+1}\binom{nk+1}{k}
=\frac{1}{(n-1)k+1}\binom{nk}{k}.
\]
The resulting series has radius
$R_n=(n-1)^{n-1}/n^n$ in $z=ba^{n-1}$. It also converges absolutely at
the boundary $|z|=R_n$. We give the derivation and the convergence argument
explicitly. The Lean code currently verifies selected coefficient facts and
algebraic identities; it does not formally verify Lagrange inversion or the
analytic convergence theorem.
\end{abstract}

\section{Scope and branch}
Throughout, $n$ is a fixed integer with $n\geq 2$, and $a,b$ are real.
The equation is a special trinomial equation, not a general polynomial of
degree $n$. It can have more than one root; the series below describes the
branch that equals $a$ at $b=0$.

If $a=0$, this branch is $x=0$, which solves $x=bx^n$. Assume henceforth
$a\neq 0$, and set
\[
y=\frac{x}{a},\qquad z=ba^{n-1},\qquad w=y-1.
\]
Then the original equation is equivalent to
\[
y=1+zy^n,\qquad w=z(1+w)^n. \tag{1}
\]
Since the derivative with respect to $w$ of
$w-z(1+w)^n$ equals $1$ at $(w,z)=(0,0)$, the analytic implicit function
theorem gives a unique analytic solution $w(z)$ near $z=0$ with $w(0)=0$.

\section{Coefficient derivation}
\begin{theorem}[Fuss--Catalan expansion]
For $|z|<R_n$, the branch in (1) is
\[
x=a\sum_{k=0}^{\infty} C_k^{(n)}z^k,\qquad
C_k^{(n)}=\frac{1}{nk+1}\binom{nk+1}{k},
\]
where $C_0^{(n)}=1$.
\end{theorem}

\begin{proof}
We use the Lagrange--B\"urmann coefficient formula. If
$W=z\phi(W)$, $\phi(0)\neq 0$, and $F$ is analytic near $0$, then for
$k\geq 1$,
\[
\operatorname{Coeff}_{z^k}\!\left(F(W(z))\right)
=\frac{1}{k}[u^{k-1}]F'(u)\phi(u)^k. \tag{2}
\]
For completeness, this formula follows by writing coefficient extraction
as a residue and making the local change of variable
$z=u/\phi(u)$:
\[
\begin{aligned}
\operatorname{Coeff}_{z^k}\!\left(F(W(z))\right)
&=\operatorname{Res}_{u=0}
\left(
\frac{F(u)\phi(u)^k}{u^{k+1}}
-\frac{F(u)\phi(u)^{k-1}\phi'(u)}{u^k}
\right)\,du\\
&=\frac{1}{k}\operatorname{Res}_{u=0}
\frac{F'(u)\phi(u)^k}{u^k}\,du\\
&=\frac{1}{k}[u^{k-1}]F'(u)\phi(u)^k.
\end{aligned}
\]
The middle equality uses that the residue of the derivative of a
Laurent series is zero.

Apply (2) to $W=w$, $\phi(u)=(1+u)^n$, and $F(u)=u$. Then
\[
[z^k]w(z)=\frac{1}{k}[u^{k-1}](1+u)^{nk}
=\frac{1}{k}\binom{nk}{k-1}.
\]
Using factorial definitions of the binomial coefficients,
\[
\frac{1}{k}\binom{nk}{k-1}
=\frac{(nk)!}{k!\,((n-1)k+1)!}
=\frac{1}{nk+1}\binom{nk+1}{k}
=C_k^{(n)}.
\]
Also $y=1+w$, so its constant coefficient is $1=C_0^{(n)}$ and its
coefficient of $z^k$ for $k\geq1$ is $C_k^{(n)}$. Since $x=ay$, the
claimed expansion follows wherever this power series converges.

Equation (1) also yields the coefficient recurrence
\[
C_0^{(n)}=1,\qquad
C_k^{(n)}
=\sum_{\substack{j_1+\cdots+j_n=k-1\\j_i\geq0}}
  C_{j_1}^{(n)}\cdots C_{j_n}^{(n)}
\quad(k\geq1), \tag{3}
\]
by comparing coefficients in $y=1+zy^n$. The convergence proof below
shows that this Taylor series is analytic on $|z|<R_n$. Both sides of
(1) are analytic there and agree near $z=0$, so the identity theorem
shows that the series satisfies the fixed-point equation throughout
$|z|<R_n$.
\end{proof}

\section{Radius of convergence and boundary}
\begin{theorem}
For fixed $n\geq2$, the power series $\sum_{k\geq0}C_k^{(n)}z^k$ has
radius of convergence
\[
R_n=\frac{(n-1)^{n-1}}{n^n}.
\]
It converges absolutely for $|z|\leq R_n$ and diverges for $|z|>R_n$.
\end{theorem}

\begin{proof}
Put $A_n=n^n/(n-1)^{n-1}$. Stirling's formula applied to
\[
C_k^{(n)}
=\frac{1}{(n-1)k+1}\binom{nk}{k}
\]
gives
\[
C_k^{(n)}
\sim
\frac{1}{n-1}\sqrt{\frac{n}{2\pi(n-1)}}\,
A_n^k k^{-3/2}. \tag{4}
\]
In particular, $\lim_{k\to\infty}(C_k^{(n)})^{1/k}=A_n$, so the
Cauchy--Hadamard formula gives radius $1/A_n=R_n$.

For $|z|<R_n$, the root test gives absolute convergence. At
$|z|=R_n$, (4) shows that the absolute values of the terms are
asymptotic to a positive constant times $k^{-3/2}$, so the series
converges absolutely by comparison with the $p$-series of exponent
$3/2$. If $|z|>R_n$, (4) shows that the terms do not tend to zero.
This proves both the radius and the stated boundary behavior.
\end{proof}

\begin{remark}
The radius is the boundary of the analytic power-series branch. The
absolute-convergence statement is not ``if and only if $|z|<R_n$'':
absolute convergence also holds at $|z|=R_n$.
\end{remark}

\section{Example}
For $n=5$,
\[
R_5=\frac{4^4}{5^5}=\frac{256}{3125}=0.08192,
\]
and the first coefficients are
\[
C_0^{(5)}=1,\quad C_1^{(5)}=1,\quad C_2^{(5)}=5,\quad
C_3^{(5)}=35,\quad C_4^{(5)}=285,\quad C_5^{(5)}=2530.
\]
For example, the parameter choices $(a,b)=(1,0.05)$ and $(2,0.002)$
give $z=0.05$ and $z=0.032$, respectively, both strictly inside the
convergence disk. These examples illustrate the theorem; numerical
residuals alone are not used as proof.

\section{Lean formalization status}
In \texttt{Basic.lean}, the coefficient is defined by the closed binomial
formula, and selected initial values and elementary residual identities
are checked. The theorem named \texttt{hassoun\_general\_root} only rearranges
the equation under the assumption that $x=a+bx^n$. Those declarations do
not prove (2), (3), the infinite-series identity, or the convergence
radius. Therefore this paper presents an analytic proof, but does not claim
that the full analytic theorem has already been machine-checked in Lean.

\section{Conclusion}
Lagrange--B\"urmann inversion gives an exact Fuss--Catalan series for the
branch through $x=a$, and Stirling's formula gives its sharp radius and
boundary convergence. The result is a standard analytic inversion
specialized to $x-bx^n=a$; any claim of novelty or full Lean verification
would require additional results beyond the derivation given here.

\end{document}
